\recipe{Butternut Squash Mac \& Cheese}
\label{recipe:butternut_squash_mac_and_cheese}
\makes{4 portions}
\worktime{\SI{30} \minute}
\totaltime{\SI{50} \minute}
\source{unknown}


\begin{ingredients}

  \item 1 medium to large butternut squash (1.5 pounds or larger), halved
  \item 1 teaspoon extra virgin olive oil
  \item 2 tablespoons unsalted butter
  \item 1 teaspoon garlic powder
  \item 1 teaspoon onion powder
  \item 1 pound short-cut pasta (pipe regate, macaroni noodles, casarecce, cavatelli, fusilli)
  \item 5 cups water
  \item 1 teaspoon fine salt, to taste
  \item 4 ounces cream cheese, cubed 
  \item 8 ounces shredded sharp cheddar cheese
  \item \qty{0.5}{\cup} finely grated
  \item Parmesan cheese, plus more for serving.
\end{ingredients}


\begin{method}
To roast the squash: Preheat the oven to 220 degrees celsuis. Line a
large, rimmed baking sheet with parchment paper for easy clean-up.
Scoop out the squash seeds with a spoon. Rub the olive oil over the cut
sides of butternut, then place them on the prepared pan, flat sides
down.

Bake until the squash flesh is easily pierced through with a fork, about
40 to 60 minutes. Once it’s cool enough to handle, peel off the skin
and discard it. Use a potato masher to mash up the squash—don’t
worry about getting it perfectly smooth. You’ll need about 2 cups
mashed squash for this recipe (you’ll likely have plenty extra, which
you can freeze for future mac and cheese). Set aside, refrigerate and/or
freeze the squash until ready to use.

To prepare the mac and cheese: Melt the butter in a large pot or Dutch
oven over medium heat. Add the garlic and onion powder and cook for
2 to 4 minutes, stirring often, until the butter is fragrant and you see
little brown flecks forming in the pan. Add the dry pasta and gently
toss to coat it in butter.

Pour in the water and salt. Add 2 cups of the mashed butternut squash.
Cover the pot and bring it to a boil over high heat. Once boiling,
remove the lid and set the timer for 8 minutes.

Cook, stirring occasionally, until the timer goes off. Do not drain the
MISSING TEXT MISSING TEXT MISSING TEXT MISSING TEXT MISSING TEXT MISSING TEXT MISSING TEXT MISSING TEXT MISSING TEXT MISSING TEXT MISSING TEXT MISSING TEXT MISSING TEXT MISSING TEXT MISSING TEXT MISSING TEXT MISSING TEXT MISSING TEXT MISSING TEXT MISSING TEXT MISSING TEXT MISSING TEXT MISSING TEXT MISSING TEXT MISSING TEXT MISSING TEXT MISSING TEXT MISSING TEXT MISSING TEXT MISSING TEXT MISSING TEXT MISSING TEXT MISSING TEXT MISSING TEXT MISSING TEXT MISSING TEXT MISSING TEXT MISSING TEXT MISSING TEXT MISSING TEXT MISSING TEXT MISSING TEXT 
  
\end{method}


\begin{figure}[h]
  	\centering
  	
  	\begin{tikzpicture}
  		% Save the image in a box to get its dimensions
  		\sbox0{\includegraphics[width = 0.6\linewidth]{images/MissingImage.png}}%
  
  		% Draw a clipping rectangle with rounded corners
  		\path[clip, rounded corners=0.5cm] (0,0) rectangle (\wd0,\ht0);
  
  		% Place the image inside the clipped area
  		\node[inner sep=0pt, outer sep=0pt, anchor=south west] at (0,0) {\usebox0};
	\end{tikzpicture}

\end{figure}


